% ICASSP 2027 submission. 4 pages technical + page 5 references only.
% DRAFT v1 -- every number traced to VERIFIED_NUMBERS.md. Nothing here is unverified.
\documentclass{article}
\usepackage{spconf,amsmath,graphicx,booktabs,microtype}

\title{Audio Language Models Encode Clinical Speech Signal They Do Not Use}

\name{Chaitanya Parwatkar, Minoo Ahmadi, Nima Kelidari, Mohammad Soleymani}
\address{University of Southern California, Los Angeles, CA, USA}

\begin{document}
\maketitle

\begin{abstract}
Audio language models are increasingly proposed for clinical speech screening, on the
assumption that their answer reflects what they heard. We show it largely does not. Across
nine audio LLMs and five clinical corpora, a linear probe on the audio encoder reaches
0.92 AUC on Parkinson's read speech while the same model's own Yes/No answer reaches 0.63.
To find out what the answer uses instead, we build a conflict set from real clinical
interviews: 195 matched segments, 74 speakers, in which the spoken words point away from
the diagnosis. Six of nine models score \emph{below chance} there (0.29--0.36) while scoring
0.70--0.93 on matched segments where words and diagnosis agree. Given only the transcript and
no audio, a model reproduces its own audio answer at $r=0.83$ and returns the identical Yes/No
on 96\% of segments. Six interventions -- twelve prompts, layer mixing, prompt-conditioned
mixing, contrastive tuning, threshold calibration and logit-space lexical subtraction -- fail
to move the answer while leaving the probe unchanged, which places the deficit in the
state-to-answer readout rather than the encoder. We release the matched conflict set.
\end{abstract}

\begin{keywords}
audio language models, clinical speech, probing, lexical shortcut, Parkinson's disease, depression
\end{keywords}

\section{Introduction}\label{sec:intro}

Speech carries measurable signs of Parkinson's disease and of depression, and audio language
models now answer clinical questions about a recording directly, in natural language. That
makes them attractive for screening: one model, one prompt, no task-specific training.

The unexamined assumption is that the answer reflects the audio. Existing clinical evaluations
report answer accuracy on data where the spoken words and the diagnosis point the same way, so
a model that reads the transcript and a model that listens to the voice are indistinguishable
by construction. For Parkinson's this matters more than usual: the disease is motor, the words
carry no diagnostic content, and a model that follows them is provably wrong. Two recent
studies reach a related conclusion on non-clinical audio \cite{listen2026,voxparadox2026}, and
multi-corpus clinical benchmarks exist \cite{speechdx2026,fend2025}, but none asks the question
inside a single audio LLM on clinical data.

We make three measurements on the same segments. (i) We compare a linear probe on the audio
encoder against the model's own emitted answer. (ii) We construct a conflict set from real
clinical interviews in which the words contradict the diagnosis, and pair every conflict
segment with a matched agreement segment. (iii) We replace the audio with its transcript under
otherwise identical conditions. All evaluation is speaker-disjoint.

The encoder carries the condition and the answer does not use it; on conflict segments the
answer follows the words; and the transcript alone reproduces the answer. Six interventions
fail to move it. We also report that one of our five corpora is separable from recording
properties alone, and that two groups of \emph{healthy} speakers recorded on different days
separate perfectly, which bounds what any result on that corpus can mean.

\section{Setup}\label{sec:setup}

\noindent\textbf{Models.} qwen2-audio, qwen2.5-omni, qwen3-omni, audio flamingo 2, audio
flamingo 3, kimi-audio, salmonn, mimo-audio and diva. All frozen, no fine-tuning.

\noindent\textbf{Corpora.} PC-GITA and Neurovoz (Spanish, Parkinson's), MDVR-KCL (English,
Parkinson's), Italian PVS (Parkinson's), DAIC-WOZ (English, depression).

\noindent\textbf{Prompt.} Verbatim, and no transcript is supplied in any audio condition:
\emph{``Based only on this recording, does this speaker show signs of X? Answer with one word,
Yes or No.''}

\noindent\textbf{Metric.} $p_{\text{yes}} = P(\text{Yes})/(P(\text{Yes})+P(\text{No}))$ read
from the logits at the first answer position, single forward pass, no generation. We also
record the pre-normalisation mass $P(\text{Yes})+P(\text{No})$, because several models place
almost none there and their AUCs are rankings rather than decisions.

\noindent\textbf{Listening window.} Model audio front-ends differ: qwen2-audio accepts 30\,s,
omni 300\,s, while DAIC interviews run 12--25 minutes. Uncontrolled, this alone produces
apparent model differences. At a matched 300\,s window flamingo 3, omni and kimi reach 0.858,
0.835 and 0.831 on depression and are statistically indistinguishable. All results below use a
window that covers the whole clip.

\noindent\textbf{Probe.} $\ell_2$-regularised logistic probe on pooled encoder states, with the
layer chosen \emph{inside} the training folds. Selecting the layer after seeing the test fold
inflates AUC by $+0.031$ on average across 18 cells, and by up to $+0.135$. Nested values move
by up to 0.049 across seeds, so we report two decimals.

\section{The gap}\label{sec:gap}

On PC-GITA read speech the encoder probe reaches \textbf{0.92} while the model's own answer on
the same clips reaches \textbf{0.63}. The information is present and the answer does not use
it. Several models do not make a decision at all: qwen2-audio answers Yes on all 1117 clips and
omni answers No on all of them, so their AUCs rank confidences that never cross the threshold.

\section{What the answer uses instead}\label{sec:conflict}

\noindent\textbf{Construction.} From the raw DAIC transcripts we take 26{,}000 participant
turns, remove interviewer questions, and merge into $\sim$30\,s segments. We keep segments
where a depressed participant says something positive or a control says something negative, and
pair each with a matched segment from the same speaker where words and diagnosis agree. This
yields 195 pairs over 74 speakers. Label and lexical direction are exactly independent
($\phi=0.000$). Unlike synthetic conflict benchmarks, nothing is generated: the voices are real
patients and the contradicting words occur naturally.

\begin{table}[t]
\centering\small
\caption{AUC against the clinical label. Below 0.5 means the model follows the words rather
than the diagnosis. No transcript is given in either arm.}
\label{tab:conflict}
\begin{tabular}{@{}lcc@{}}
\toprule
model & conflict & agreement \\
\midrule
salmonn            & \textbf{0.29} & 0.85 \\
diva               & \textbf{0.30} & 0.70 \\
qwen2.5-omni       & \textbf{0.31} & 0.85 \\
qwen3-omni         & \textbf{0.35} & 0.87 \\
qwen2-audio        & \textbf{0.35} & 0.92 \\
audio flamingo 3   & \textbf{0.36} & 0.93 \\
mimo-audio         & 0.47 & 0.82 \\
kimi-audio         & 0.52 & 0.87 \\
audio flamingo 2   & 0.55 & 0.57 \\
\bottomrule
\end{tabular}
\end{table}

Six of nine models score below chance (Table~\ref{tab:conflict}), across five independent
labs. Audio flamingo 2 is flat on both arms and therefore has no signal to invert.

\noindent\textbf{The transcript reproduces the answer.} Given only the Whisper transcript and
no audio, under the same prompt, a model's answer correlates with its own audio answer at
$r=0.83$ and is the identical Yes/No on 96\% of segments. Across the full grid of 30
(model, corpus) cells, \emph{no model beats its own transcript on depression} (0 of 6; median
text 0.81 vs audio 0.83). The 11 cells where audio wins are Parkinson's read tasks where the
transcript is already near chance.

\section{Six interventions, and what their failure shows}\label{sec:interventions}

Twelve prompt wordings never lift conflict AUC above chance. Learned layer mixing reaches 0.909
against 0.916 for honestly choosing one layer; conditioning the mixture on the prompt adds
0.0004. Contrastive tuning and threshold calibration move the decision boundary without
changing what is discriminated.

\noindent\textbf{Lexical subtraction.} Two forward passes, one on audio and one on the
transcript alone, subtracting text logits from audio logits at weight $\alpha$. At
$\alpha=0.9$ the model stops following the words (agreement with the lexical direction falls
$0.70\rightarrow0.50$) but conflict AUC only reaches 0.502; beyond that it overshoots into
inversion. Fitted on held-out speakers, $\alpha=0$ is chosen in four of five folds.

\noindent\textbf{Encoder steering.} We add a multiple of the difference-in-means diagnosis
direction to the encoder output before the projector, keeping projector and LLM frozen. The
agreement arm is already significantly damaged at the smallest coefficient while conflict has
not moved, and a \emph{random direction of the same norm displaces the answer equally}. Across
20 steered conditions, conflict AUC and lexical AUC correlate at $-0.916$: the apparent gain is
the lexical signal being destroyed, not the voice emerging.

\noindent\textbf{A published fix does not transfer.} PCLM+DPO \cite{voxparadox2026}, which
lifts synthetic paralinguistic accuracy substantially, leaves conflict AUC at 0.365 against a
base of 0.350, and every clinical delta falls inside the base model's bootstrap interval. The
prompt-conditioned gate is close to a no-op on clinical prompts, placing 66\% of its mass on
the final encoder layer against 37\% on VoxParadox-style prompts.

Read together these are not six unrelated failures. Every intervention that leaves the encoder
intact fails to move the answer, while the probe on that same encoder is unchanged. The deficit
is in the state-to-answer readout, not in the representation.

\section{Corpus validity}\label{sec:validity}

Six recording measurements containing no speech content separate patients from controls at
0.995 on Italian PVS; the silent portions alone reach 1.0, and the file container alone,
with no audio at all, reaches 0.776. The nested probe selects layer 0 in every fold and seed.
Decisively, \emph{two groups of healthy elderly controls from the same cohort, recorded on
different days, separate at speaker-level AUC 1.000} ($p=0.039$), while other pairs of elderly
sessions do not separate at all (0.11--0.82, n.s.). Published accuracies of 98--99\% on this
corpus are therefore not evidence about pathology. We report Italian as a negative control and
exclude it from all model claims. The other four corpora have floors of 0.45--0.80 from
recording properties alone and retain signal above them.

\section{Limitations}\label{sec:limits}

The conflict set is depression only; Parkinson's tasks are fixed read sentences and sustained
vowels, so no naturally contradicting words exist there. MDVR-KCL has 37 speakers and its
intervals are correspondingly wide. Word-following is associated with whether the audio encoder
was trained on a transcription objective (6 of 7 ASR-trained encoders invert vs 0 of 2
non-ASR), but with nine models this is suggestive rather than significant (Fisher exact
$p=0.08$), the predicate was chosen after seeing the results, and kimi-audio is a
counterexample. Depression scores track PTSD severity as strongly as depression severity
(0.602 vs 0.585), so the models detect distress rather than depression specifically.

\section{Conclusion}\label{sec:conclusion}

Across nine audio LLMs the clinical signal is present in the encoder and absent from the
answer, and what the answer uses instead is the transcript. Six interventions that leave the
encoder intact fail to change this. Evaluations that score these models only where words and
diagnosis agree cannot distinguish a model that listens from one that reads.

\bibliographystyle{IEEEbib}
\bibliography{refs}

\end{document}
